% Part: first-order-logic
% Chapter: tableaux
% Section: provability-consistency

\documentclass[../../../include/open-logic-section]{subfiles}

\begin{document}

\iftag{FOL}
      {\olfileid{fol}{tab}{prv}}
      {\olfileid{pl}{tab}{prv}}

\olsection{\usetoken{S}{derivability} અને સુસંગતતા}

હવે આપણે !!{derivability} સંબંધના કેટલાક ગુણધર્મો સ્થાપિત કરીશું. આ
ગુણધર્મો સ્વતંત્ર રીતે રસપ્રદ છે, પરંતુ દરેક પૂર્ણતાના પ્રમેયની
સાબિતીમાં ભૂમિકા ભજવશે.

\begin{prop}\ollabel{prop:provability-contr}
  જો $\Gamma \Proves !A$ હોય અને $\Gamma \cup \{!A\}$ અસુસંગત હોય,
  તો $\Gamma$ અસુસંગત છે.
\end{prop}

\begin{proof}
એવા સાન્ત ગણ $\Gamma_0 = \{!B_1, \dots, !B_n\}$ અને $\Gamma_1
  =\{!C_1, \dots, !C_m\} \subseteq \Gamma$ છે કે
  \begin{align*}
    \{\sFmla{\False}{!A}, &
    \sFmla{\True}{!B_1}, \dots, \sFmla{\True}{!B_n}\} \\
    \{\sFmla{\True}{!A}, &
    \sFmla{\True}{!C_1}, \dots, \sFmla{\True}{!C_m}\}
  \end{align*}
  માટે બંધ !!{tableau}s છે. $!A$ પર \Cut{} નિયમ વાપરીને આપણે તેમને એક
  એવા બંધ !!{tableau}માં જોડી શકીએ છીએ જે દર્શાવે છે કે
  $\Gamma_0 \cup \Gamma_1$ અસુસંગત છે. $\Gamma_0 \subseteq \Gamma$ અને
  $\Gamma_1 \subseteq \Gamma$ હોવાથી $\Gamma_0 \cup \Gamma_1
  \subseteq \Gamma$; તેથી $\Gamma$ અસુસંગત છે.
  \sourcecorrection{OLTAB-004}{મૂળની
  \texttt{provability-consistency.tex}, પંક્તિ~26માં બીજા સાન્ત ગણને
  $\{!C_1,\dots,!C_n\}$ લખ્યો છે, પરંતુ તરત પછીનું પ્રદર્શન અને બાકીની
  સાબિતી તેનો અંતિમ સૂચકાંક~$m$ રાખે છે. અહીં $\{!C_1,\dots,!C_m\}$
  પુનઃસ્થાપિત કર્યો છે.}
\end{proof}

\begin{prop}
\ollabel{prop:prov-incons}
$\Gamma \Proves !A$ જો અને માત્ર જો $\Gamma \cup \{\lnot !A\}$
અસુસંગત છે.
\end{prop}

\begin{proof}
પહેલાં ધારો કે $\Gamma \Proves !A$, એટલે કે નીચેના ગણ માટે બંધ
!!{tableau} છે:
\[
\{\sFmla{\False}{!A},
\sFmla{\True}{!B_1}, \dots, \sFmla{\True}{!B_n}\}
\]
$\TRule{\True}{\lnot}$ નિયમ વાપરીને તેને નીચેના ગણ માટેના બંધ
!!{tableau}માં ફેરવી શકાય છે:
\[
\{\sFmla{\True}{\lnot !A},
\sFmla{\True}{!B_1}, \dots, \sFmla{\True}{!B_n}\}.
\]

બીજી બાજુ, જો પછીના ગણ માટે બંધ !!{tableau} હોય, તો પહેલી
ધારણા~$\sFmla{\True}{\lnot !A}$ અને એ ધારણાને
\TRule{\True}{\lnot} લાગુ કરવાથી મળતું દરેક સૂત્ર દૂર કરીને, તથા ધારણા
$\sFmla{\False}{!A}$ ઉમેરીને તેને પહેલા ગણના બંધ !!{tableau}માં ફેરવી
શકીએ છીએ. ખરેખર, કોઈ શાખા પહેલાં $\sFmla{\True}{\lnot !A}$ને
\TRule{\True}{\lnot} લાગુ કરવાના ફલિત $\sFmla{\False}{!A}$ને કારણે બંધ
હોય, તો નવા !!{tableau}માં તેને અનુરૂપ શાખા પણ બંધ છે. જૂના ટેબ્લોમાં
કોઈ શાખા ધારણા $\sFmla{\True}{\lnot !A}$ની સાથે
$\sFmla{\False}{\lnot !A}$ ધરાવવાને કારણે બંધ હોય, તો
$\sFmla{\False}{\lnot !A}$ને $\TRule{\False}{\lnot}$ લાગુ કરીને
$\sFmla{\True}{!A}$ મેળવી તેને બંધ શાખામાં ફેરવી શકીએ છીએ. આપણે
$\sFmla{\False}{!A}$ને ધારણા તરીકે ઉમેર્યું હોવાથી આ શાખા બંધ થાય છે.
\end{proof}

\begin{prob}
$\Gamma \Proves \lnot !A$ હોય જો અને માત્ર જો
$\Gamma \cup \{!A\}$ અસુસંગત હોય, એ સાબિત કરો.
\end{prob}

\begin{prop}\ollabel{prop:explicit-inc}
  જો $\Gamma \Proves !A$ અને $\lnot !A \in \Gamma$ હોય, તો $\Gamma$
  અસુસંગત છે.
\end{prop}

\begin{proof}
  ધારો કે $\Gamma \Proves !A$ અને $\lnot !A \in \Gamma$. ત્યારે એવાં
  $!B_1$, \dots, $!B_n \in \Gamma$ છે કે
  \[
  \{\sFmla{\False}{!A}, \sFmla{\True}{!B_1}, \dots, \sFmla{\True}{!B_n}\}
  \]
  માટે બંધ !!{tableau} છે. ધારણા \sFmla{\False}{!A}ને
  \sFmla{\True}{\lnot !A}થી બદલો અને \sFmla{\True}{\lnot !A}ને
  \TRule{\True}{\lnot} લાગુ કરવાનું ફલિત ધારણાઓ પછી દાખલ કરો. જૂના
  !!{tableau}માં પંક્તિ~$1$ના આધારે સમર્થિત દરેક !!{sentence} હવે
  પંક્તિ~$n+1$ના આધારે સમર્થિત છે. તેથી જૂનો !!{tableau} બંધ હોય તો
  નવો પણ બંધ છે. બધી ધારણાઓ~$\Gamma$માં હોવાથી તે દર્શાવે છે કે
  $\Gamma$ અસુસંગત છે.
  \sourcecorrection{OLTAB-005}{મૂળની
  \texttt{provability-consistency.tex}, પંક્તિ~93માં
  $\TRule{\True}{\lnot}$ને $\sFmla{\False}{!A}$ પર લાગુ કરવાનું લખ્યું
  છે. નિયમ અને પૂર્વવાક્ય મુજબ તેનું આધારવાક્ય
  $\sFmla{\True}{\lnot !A}$ છે; અહીં એ જ પુનઃસ્થાપિત કર્યું છે.}
\end{proof}

\begin{prop}\ollabel{prop:provability-exhaustive}
  જો $\Gamma \cup \{!A\}$ અને $\Gamma \cup \{\lnot !A\}$ બંને
  અસુસંગત હોય, તો $\Gamma$ અસુસંગત છે.
\end{prop}

\begin{proof}
  જો એવાં $!B_1$, \dots, $!B_n \in \Gamma$ અને $!C_1$, \dots,
  $!C_m \in \Gamma$ હોય કે
  \begin{align*}
    \{\sFmla{\True}{!A}, &
    \sFmla{\True}{!B_1}, \dots, \sFmla{\True}{!B_n}\} \text{ અને}\\
    \{\sFmla{\True}{\lnot !A}, &
    \sFmla{\True}{!C_1}, \dots, \sFmla{\True}{!C_m}\}
  \end{align*}
  બંને માટે બંધ !!{tableau}s હોય, તો
  $\sFmla{\True}{!B_1}$, \dots, $\sFmla{\True}{!B_n}$ની સાથે
  $\sFmla{\True}{!C_1}$, \dots, $\sFmla{\True}{!C_m}$ને ધારણાઓ તરીકે
  લઈને અને પછી \Cut{} નિયમ લાગુ કરીને એક જ સંયુક્ત !!{tableau} રચી
  શકીએ છીએ, જે દર્શાવે છે કે $\Gamma$ અસુસંગત છે. તેનાથી બે શાખાઓ મળે
  છે: એક $\sFmla{\True}{!A}$થી અને બીજી $\sFmla{\False}{!A}$થી શરૂ થાય
  છે.

  ડાબી બાજુએ પહેલા !!{tableau}નો તેની ધારણાઓની નીચેનો ભાગ ઉમેરો. અહીં
  દરેક નિયમનો ઉપયોગ હજી યોગ્ય છે, કારણ કે પહેલા !!{tableau}ની
  $\sFmla{\True}{!A}$ સહિતની દરેક ધારણા ઉપલબ્ધ છે. તેથી
  $\sFmla{\True}{!A}$ની નીચેની દરેક શાખા બંધ થાય છે.

  જમણી બાજુએ બીજા !!{tableau}નો તેની ધારણાઓની નીચેનો ભાગ ઉમેરો, પરંતુ
  $\sFmla{\True}{\lnot !A}$ને~$\TRule{\True}{\lnot}$ લાગુ કરવાનાં બધાં
  પરિણામો દૂર કરો. $\sFmla{\True}{\lnot !A}$ને
  $\TRule{\True}{\lnot}$ લાગુ કરવાનું ફલિત $\sFmla{\False}{!A}$ છે; તે
  છતાં ઉપલબ્ધ છે, કારણ કે તે સંયુક્ત !!{tableau}ની જમણી બાજુના \Cut{}
  નિયમનું ફલિત છે.

  બીજા ટેબ્લોમાં કોઈ શાખા ધારણા $\sFmla{\True}{\lnot !A}$ (જે હવે
  સંયુક્ત !!{tableau}માં ધારણા તરીકે આવતી નથી) અને
  $\sFmla{\False}{\lnot !A}$ બંને ધરાવવાને કારણે બંધ થઈ હોય, તો
  $\sFmla{\False}{\lnot !A}$ને $\TRule{\False}{\lnot}$ લાગુ કરીને
  $\sFmla{\True}{!A}$ મેળવી શકીએ છીએ. હવે સંયુક્ત !!{tableau}માં તેને
  અનુરૂપ શાખા પણ બંધ થાય છે, કારણ કે તેમાં \Cut{} નિયમનું જમણી બાજુનું
  ફલિત $\sFmla{\False}{!A}$ છે. બીજા !!{tableau}ની કોઈ શાખા અન્ય કોઈ
  કારણે બંધ થઈ હોય, તો સંયુક્ત !!{tableau}માં તેને અનુરૂપ શાખા પણ બંધ
  થાય છે, કારણ કે જૂના બીજા !!{tableau}ની શાખામાં આવતાં
  $\sFmla{\True}{\lnot !A}$ સિવાયનાં બધાં !!{signed formula}s સંયુક્ત
  !!{tableau}ની તેને અનુરૂપ શાખામાં પણ આવે છે.
  \end{proof}

\end{document}
